% from ega2-1-fr.tex
\oldpage[II]{145}
\item[\rm c)] Pour tout $Y$-schéma de la forme $Y'=\operatorname{Spec}(A)$,
où $A$ est un anneau de valuation (resp. un anneau de valuation discrète)
l'application canonique (\hyperref[II.7.3.2.2-fr]{7.3.2.2}) relative au
$Y'$-préschéma $X_{(Y')}$ est surjective (resp. bijective).
\end{enumerate}
\end{theorem}

L'équivalence de b) et c) résulte aussitôt de
(\hyperref[I.3.3.14-fr]{I, 3.3.14})~; a) implique c), car a) entraîne, dans
l'une ou l'autre hypothèse, que $f_{(Y')}$ est séparé
(\hyperref[I.5.5.1-fr]{I, 5.5.1, (iv)}) et fermé, et il suffit d'appliquer
(\hyperref[II.7.3.3-fr]{7.3.3}). Reste à prouver que b) entraîne a).
Plaçons-nous d'abord dans le cas où $Y$ est quelconque, $f$ séparé et
quasi-compact. Si la condition b) est vérifiée par $f$, elle l'est aussi par
$f_{(Y'')}:X_{(Y'')}\to Y''$, où $Y''$ est un $Y$-préschéma quelconque, en
vertu de l'équivalence de b) et c), et du fait que
$X_{(Y'')}\times_{Y''}Y'=X\times_YY'$ pour tout morphisme $Y'\to Y''$
(\hyperref[I.3.3.9.1-fr]{I, 3.3.9.1})~; comme en outre $f_{(Y'')}$ est séparé
et quasi-compact quand $f$ l'est
(\hyperref[I.5.5.1-fr]{I, 5.5.1, (iv)} et
\hyperref[I.6.6.4-fr]{6.6.4, (iii)}), on est ramené à prouver que b) entraîne
que $f$ est \emph{fermé}. Pour cela, il suffit de vérifier la condition b) de
(\hyperref[II.7.2.1-fr]{7.2.1}). Soient donc $x\in X$, $y'$ une spécialisation
de $y=f(x)$, distincte de $y$~; en vertu de
(\hyperref[II.7.1.4-fr]{7.1.4}), il y a un schéma $Y'$, spectre d'un anneau de
valuation, et un morphisme séparé $g:Y'\to Y$ tels que, si $a$ désigne le point
fermé et $b$ le point générique de $Y'$, on ait $g(a)=y'$, $g(b)=y$, et qu'il
existe un $k(y)$-homomorphisme $k(x)\to k(b)$. Ce dernier correspond
canoniquement à un $Y$-morphisme $\operatorname{Spec}(k(b))\to X$
(\hyperref[I.2.4.6-fr]{I, 2.4.6}), et il résulte donc de b) qu'il existe un
$Y$-morphisme $h:Y'\to X$ auquel correspond le morphisme précédent. On a
alors $h(b)=x$~; si l'on pose $h(a)=x'$, $x'$ est spécialisation de $x$, et
l'on a $f(x')=f(h(a))=g(a)=y'$.

Si maintenant $Y$ est localement noethérien et $f$ de type fini, l'hypothèse
b) entraîne d'abord que $f$ est \emph{séparé}, en vertu de
(\hyperref[II.7.2.3-fr]{7.2.3}), le morphisme diagonal $X\to X\times_YX$
étant quasi-compact (\hyperref[II.7.2.4-fr]{7.2.4}). En outre, pour vérifier
que $f$ est propre, il suffit de montrer que $f_{(Y'')}:X_{(Y'')}\to Y''$ est
\emph{fermé} pour tout $Y$-préschéma $Y''$ \emph{de type fini}, compte tenu de
(\hyperref[II.5.6.3-fr]{5.6.3}). Comme alors $Y''$ est localement noethérien,
on peut reprendre le raisonnement fait dans le premier cas en prenant pour
$Y'$ un spectre d'anneau de valuation discrète, et en appliquant
(\hyperref[II.7.1.9-fr]{7.1.9}) au lieu de
(\hyperref[II.7.1.4-fr]{7.1.4}).

\begin{remark}[7.3.9]
\label{II.7.3.9-fr}
\begin{enumerate}
\item[\rm (i)] Lorsque $Y$ est un préschéma quelconque et $f$ un morphisme
séparé, pour que $f$ soit universellement fermé, il \emph{suffit} que la
condition b) ou la condition c) soit vérifiée pour les anneaux de valuation
$A$ \emph{complets} et dont le corps des restes est \emph{algébriquement
clos}~; cela résulte en effet de la démonstration et de
(\hyperref[II.7.1.4-fr]{7.1.4}).
\item[\rm (ii)] On déduit du critère c) de
(\hyperref[II.7.3.8-fr]{7.3.8}) une nouvelle démonstration du fait qu'un
morphisme projectif $X\to Y$ est fermé (\hyperref[II.5.5.3-fr]{5.5.3}), plus
proche des méthodes classiques. On peut en effet supposer $Y$ affine, et par
suite $X$ identifié à un sous-préschéma fermé d'un fibré projectif
$\mathbf{P}_Y^n$ (\hyperref[II.5.3.3-fr]{5.3.3})~; pour prouver que $X\to Y$
est fermé, il suffit de vérifier qu'il en est ainsi du morphisme structural
$\mathbf{P}_Y^n\to Y$, et le critère c) de
(\hyperref[II.7.3.8-fr]{7.3.8}), joint à
(\hyperref[II.4.1.3.1-fr]{4.1.3.1}), prouve qu'on est ramené à démontrer le
fait suivant~: \emph{si $Y$ est le spectre d'un anneau de valuation $A$, de
corps des fractions $K$, tout point de $\mathbf{P}_Y^n$ à valeurs dans $K$
provient (par restriction au point générique de $Y$) d'un point de
$\mathbf{P}_Y^n$ à valeurs dans $A$}. Or, tout $\mathcal{O}_Y$-Module
inversible est trivial (\hyperref[I.2.4.8-fr]{I, 2.4.8})~; donc, il résulte de
(\hyperref[II.4.2.6-fr]{4.2.6}) qu'un point de $\mathbf{P}_Y^n$ à valeurs dans
$K$ s'identifie à une classe d'éléments $(\zeta c_0,\zeta c_1,\ldots,\zeta
c_n)$ de $K$, où $\zeta\neq0$ et les $c_i$ sont des éléments non tous nuls de
$K$. Or, en multipliant les $c_i$ par un élément de $A$ de
